% GENERATED FILE. Edit canonical lessons, then run npm run generate:books.
% canonical-source-hash: fnv1a64:65af0e1c34c87eac
% canonical-lessons: TE-C455-paryatakudu, TE-C455-paryatana, TE-C455-gaidu, TE-C455-pravesam, TE-C455-viharayatra

\chapter{Being a Tourist}
\label{ch:te-being-a-tourist}
\hlchaptermodality{455}

\begin{chapteropening}
\textbf{By the end of this chapter:} \emph{I can say a tourist, a tour, a tour guide, entry and a pleasure trip.}

Complete the last lesson of chapter 455: I can say a tourist, a tour, a tour guide, entry and a pleasure trip.
\end{chapteropening}

\section[\texorpdfstring{\te{పర్యాటకుడు} (paryāṭakuḍu)}{పర్యాటకుడు (paryāṭakuḍu)}]{\te{పర్యాటకుడు} (paryāṭakuḍu) — a tourist}

\label{lesson:TE-C455-paryatakudu}



\begin{quote}
Before the new one: say the Telugu for transport, then the Telugu for to see off.
\end{quote}

\subsection*{You'll want to know: \te{పర్యాటకుడు}}
\textbf{\te{పర్యాటకుడు}} — \emph{paryāṭakuḍu} — \textquotedblleft{}a tourist\textquotedblright{}.

\begin{quote}
\te{వేసవిలో} \te{ఆ} \te{కొండ} \te{మీద} \te{పర్యాటకులు} \te{ఎక్కువగా} \te{ఉంటారు}.
\end{quote}

\begin{quote}
\emph{vēsavilō ā koṇḍa mīda paryāṭakulu ekkuvagā uṇṭāru.}
\end{quote}

\begin{quote}
In summer there are a lot of tourists on that hill.
\end{quote}

Plural: \te{పర్యాటకులు}, \emph{paryāṭakulu}.

Plural \te{పర్యాటకులు}. \te{పర్యాటక} \te{ప్రదేశం} is a tourist spot.

\begin{cousinweb}
\textbf{\te{పర్యాటకుడు}} (\emph{paryāṭakuḍu}) is borrowed from Sanskrit \emph{paryaṭaka}, which Monier-Williams glosses as a wanderer, a vagabond, from \emph{pari-aṭ}, \textquotedblleft{}to roam about\textquotedblright{}. Telugu gave the roamer a camera.
\end{cousinweb}

\begin{grammarlens}[title={What you've built}]
One more word for getting around.
\end{grammarlens}

\subsection*{Your turn}
\begin{itemize}
\raggedright
  \item \emph{Say it:} \emph{paryāṭakuḍu}
  \item \emph{Say it:} \emph{paryāṭakuḍu}, once more
  \item \emph{Say it:} \emph{vēsavilō ā koṇḍa mīda paryāṭakulu ekkuvagā uṇṭāru.} — In summer there are a lot of tourists on that hill.
  \item \emph{Recall:} say the Telugu for transport, then the Telugu for to see off, then say \emph{paryāṭakuḍu} again
\end{itemize}

\begin{tcolorbox}[breakable,colback=teal!4,colframe=teal!35!black,title={Before you move on}]
What does \te{పర్యాటకుడు} mean? (\textquotedblleft{}A tourist\textquotedblright{}.) Say it once more.
\end{tcolorbox}
\section[\texorpdfstring{\te{పర్యటన} (paryaṭana)}{పర్యటన (paryaṭana)}]{\te{పర్యటన} (paryaṭana) — a tour; an official visit}

\label{lesson:TE-C455-paryatana}



\begin{quote}
Before the new one: say the Telugu for to see off, then the Telugu for a tourist.
\end{quote}

\subsection*{You'll want to know: \te{పర్యటన}}
\textbf{\te{పర్యటన}} — \emph{paryaṭana} — \textquotedblleft{}a tour; an official visit\textquotedblright{}.

\begin{quote}
\te{మంత్రి} \te{పర్యటన} \te{వల్ల} \te{రోడ్లు} \te{మూసివేశారు}.
\end{quote}

\begin{quote}
\emph{mantri paryaṭana valla rōḍlu mūsivēśāru.}
\end{quote}

\begin{quote}
The roads were closed because of the minister's visit.
\end{quote}

A tour of places, and also the official visit of a leader. \te{పర్యటించు}, to tour, is its verb.

\begin{cousinweb}
\textbf{\te{పర్యటన}} (\emph{paryaṭana}) is borrowed from Sanskrit \emph{paryaṭana}, \textquotedblleft{}wandering about, roaming through\textquotedblright{} (Monier-Williams).
\end{cousinweb}

\begin{grammarlens}[title={What you've built}]
One more word for getting around.
\end{grammarlens}

\subsection*{Your turn}
\begin{itemize}
\raggedright
  \item \emph{Say it:} \emph{paryaṭana}
  \item \emph{Say it:} \emph{paryaṭana}, once more
  \item \emph{Say it:} \emph{mantri paryaṭana valla rōḍlu mūsivēśāru.} — The roads were closed because of the minister's visit.
  \item \emph{Recall:} say the Telugu for to see off, then the Telugu for a tourist, then say \emph{paryaṭana} again
\end{itemize}

\begin{tcolorbox}[breakable,colback=teal!4,colframe=teal!35!black,title={Before you move on}]
What does \te{పర్యటన} mean? (\textquotedblleft{}A tour; an official visit\textquotedblright{}.) Say it once more.
\end{tcolorbox}
\section[\texorpdfstring{\te{గైడు} (gaiḍu)}{గైడు (gaiḍu)}]{\te{గైడు} (gaiḍu) — a tour guide}

\label{lesson:TE-C455-gaidu}



\begin{quote}
Before the new one: say the Telugu for a tourist, then the Telugu for a tour.
\end{quote}

\subsection*{You'll want to know: \te{గైడు}}
\textbf{\te{గైడు}} — \emph{gaiḍu} — \textquotedblleft{}a tour guide\textquotedblright{}.

\begin{quote}
\te{గైడు} \te{మాకు} \te{కోట} \te{చరిత్ర} \te{చెప్పాడు}.
\end{quote}

\begin{quote}
\emph{gaiḍu māku kōṭa caritra ceppāḍu.}
\end{quote}

\begin{quote}
The guide told us the history of the fort.
\end{quote}

It takes -\te{ఉ} like other English loans. A guidebook is also called a \te{గైడు}.

\begin{cousinweb}
\textbf{\te{గైడు}} (\emph{gaiḍu}) is a loan from English \emph{guide}, with the Telugu final -\te{ఉ}.
\end{cousinweb}

\begin{grammarlens}[title={What you've built}]
One more word for getting around.
\end{grammarlens}

\subsection*{Your turn}
\begin{itemize}
\raggedright
  \item \emph{Say it:} \emph{gaiḍu}
  \item \emph{Say it:} \emph{gaiḍu}, once more
  \item \emph{Say it:} \emph{gaiḍu māku kōṭa caritra ceppāḍu.} — The guide told us the history of the fort.
  \item \emph{Recall:} say the Telugu for a tourist, then the Telugu for a tour, then say \emph{gaiḍu} again
\end{itemize}

\begin{tcolorbox}[breakable,colback=teal!4,colframe=teal!35!black,title={Before you move on}]
What does \te{గైడు} mean? (\textquotedblleft{}A tour guide\textquotedblright{}.) Say it once more.
\end{tcolorbox}
\section[\texorpdfstring{\te{ప్రవేశం} (pravēśaṁ)}{ప్రవేశం (pravēśaṁ)}]{\te{ప్రవేశం} (pravēśaṁ) — entry; admission}

\label{lesson:TE-C455-pravesam}



\begin{quote}
Before the new one: say the Telugu for a tour, then the Telugu for a tour guide.
\end{quote}

\subsection*{You'll want to know: \te{ప్రవేశం}}
\textbf{\te{ప్రవేశం}} — \emph{pravēśaṁ} — \textquotedblleft{}entry; admission\textquotedblright{}.

\begin{quote}
\te{మ్యూజియంలో} \te{పిల్లలకు} \te{ప్రవేశం} \te{ఉచితం}.
\end{quote}

\begin{quote}
\emph{myūjiyaṁlō pillalaku pravēśaṁ ucitaṁ.}
\end{quote}

\begin{quote}
Entry to the museum is free for children.
\end{quote}

\te{ప్రవేశం} \te{లేదు}, no entry. \te{ప్రవేశించు}, to enter, is its verb.

\begin{cousinweb}
\textbf{\te{ప్రవేశం}} (\emph{pravēśaṁ}) is borrowed from Sanskrit \emph{praveśa}, \textquotedblleft{}entering, entrance\textquotedblright{} (Monier-Williams), from \emph{pra-viś}, \textquotedblleft{}to enter\textquotedblright{}.
\end{cousinweb}

\begin{grammarlens}[title={What you've built}]
One more word for getting around.
\end{grammarlens}

\subsection*{Your turn}
\begin{itemize}
\raggedright
  \item \emph{Say it:} \emph{pravēśaṁ}
  \item \emph{Say it:} \emph{pravēśaṁ}, once more
  \item \emph{Say it:} \emph{myūjiyaṁlō pillalaku pravēśaṁ ucitaṁ.} — Entry to the museum is free for children.
  \item \emph{Recall:} say the Telugu for a tour, then the Telugu for a tour guide, then say \emph{pravēśaṁ} again
\end{itemize}

\begin{tcolorbox}[breakable,colback=teal!4,colframe=teal!35!black,title={Before you move on}]
What does \te{ప్రవేశం} mean? (\textquotedblleft{}Entry; admission\textquotedblright{}.) Say it once more.
\end{tcolorbox}
\section[\texorpdfstring{\te{విహారయాత్ర} (vihārayātra)}{విహారయాత్ర (vihārayātra)}]{\te{విహారయాత్ర} (vihārayātra) — a pleasure trip; an excursion}

\label{lesson:TE-C455-viharayatra}



\begin{quote}
Before the new one: say the Telugu for a tour guide, then the Telugu for entry.
\end{quote}

\subsection*{You'll want to know: \te{విహారయాత్ర}}
\textbf{\te{విహారయాత్ర}} — \emph{vihārayātra} — \textquotedblleft{}a pleasure trip; an excursion\textquotedblright{}.

\begin{quote}
\te{బడి} \te{పిల్లలు} \te{అడవికి} \te{విహారయాత్రకు} \te{వెళ్ళారు}.
\end{quote}

\begin{quote}
\emph{baḍi pillalu aḍaviki vihārayātraku veḷḷāru.}
\end{quote}

\begin{quote}
The schoolchildren went on an excursion to the forest.
\end{quote}

\te{విహార}-, roaming for pleasure, + \te{యాత్ర}, a journey. A school trip is a \te{విహారయాత్ర}.

\begin{cousinweb}
\textbf{\te{విహారయాత్ర}} (\emph{vihārayātra}) is a compound of two words borrowed from Sanskrit: \emph{vihāra}, \textquotedblleft{}roaming, a pleasure ground\textquotedblright{} (CDIAL 12033), and \emph{yātrā}, \textquotedblleft{}a journey\textquotedblright{} (Monier-Williams).
\end{cousinweb}

\begin{grammarlens}[title={What you've built}]
One more word for getting around.
\end{grammarlens}

\subsection*{Your turn}
\begin{itemize}
\raggedright
  \item \emph{Say it:} \emph{vihārayātra}
  \item \emph{Say it:} \emph{vihārayātra}, once more
  \item \emph{Say it:} \emph{baḍi pillalu aḍaviki vihārayātraku veḷḷāru.} — The schoolchildren went on an excursion to the forest.
  \item \emph{Recall:} say the Telugu for a tour guide, then the Telugu for entry, then say \emph{vihārayātra} again
\end{itemize}

\begin{tcolorbox}[breakable,colback=teal!4,colframe=teal!35!black,title={Before you move on}]
What does \te{విహారయాత్ర} mean? (\textquotedblleft{}A pleasure trip; an excursion\textquotedblright{}.) Say it once more.
\end{tcolorbox}
